% !TeX root = ../偏微分方程笔记.tex
% 由主文档通过 \input 载入；请编译主文档。
\setcounter{chapter}{9} % 保留原书章号。
\chapter{拉普拉斯方程与调和函数}
\label{chap:harmonic}

平衡场没有独立的时间演化，但其内部值仍受强烈约束：一个点的值等于周围球面的平均。平均值性质是极值原理和内部平滑性的共同来源。

\section{调和函数与格林恒等式}
\begin{definition}
若 $u\in C^2(\Omega)$ 且 $\Delta u=0$，则称 $u$ 为\emph{调和函数（harmonic function）}。
\end{definition}
调和不意味着所有二阶导数都为零，而意味着它们的和为零。例如 $u(x,y)=x^2-y^2$ 满足 $u_{xx}=2$、$u_{yy}=-2$，两者抵消，所以调和；$x^2+y^2$ 的拉普拉斯算子为 $4$，则不是调和函数。前者在原点是鞍点，与“非常数调和函数没有内部最大值或最小值”相容。线性函数也是调和函数，但两个调和函数的乘积一般不调和，因为乘积公式会多出 $2\nabla u\cdot\nabla v$。

设 $\Omega$ 为有界 $C^1$ 区域，$u,v\in C^2(\overline\Omega)$。对 $u\nabla v$ 使用散度定理得到\emph{格林第一恒等式（Green's first identity）}
\begin{equation}
 \int_\Omega(\nabla u\cdot\nabla v+u\Delta v)\,dx
 =\int_{\partial\Omega}u\partial_\nu v\,dS.
 \label{eq:green-first}
\end{equation}
交换 $u,v$ 并相减，得到格林第二恒等式
\begin{equation}
 \int_\Omega(u\Delta v-v\Delta u)\,dx
 =\int_{\partial\Omega}(u\partial_\nu v-v\partial_\nu u)\,dS.
 \label{eq:green-second}
\end{equation}
本章使用外法向；改变方向会同时改变全部边界项的符号。

\section{平均值性质与内部正则性}
记 $\sigma_{n-1}$ 为单位球面面积。若 $\overline{B_R(x_0)}\subset\Omega$，球面平均 $m(r)$ 的导数满足
\[
 m'(r)=\frac1{\sigma_{n-1}r^{n-1}}\int_{\partial B_r(x_0)}\partial_\nu u\,dS
 =\frac1{\sigma_{n-1}r^{n-1}}\int_{B_r(x_0)}\Delta u\,dx=0.
\]
令 $r\downarrow0$，再对半径积分，得到
\begin{equation}
 u(x_0)=\frac1{|\partial B_r|}\int_{\partial B_r(x_0)}u\,dS
       =\frac1{|B_r|}\int_{B_r(x_0)}u\,dx.
 \label{eq:harmonic-mean}
\end{equation}
这称\emph{平均值性质（mean value property）}。

若连续函数在所有内含小球上满足平均值性质，则与任意径向光滑单位质量磨光核卷积后在内部仍等于自身，故它内部光滑。对球平均作二阶泰勒展开再令半径趋于零，便得 $\Delta u=0$。同一卷积表示也说明调和函数的高阶导数可由稍大邻域中的函数值控制；内部光滑并不保证连续到任意形状的边界。

\section{最大值原理与唯一性}
\begin{theorem}
设 $\Omega$ 有界、连通，$u\in C^2(\Omega)\cap C(\overline\Omega)$ 调和，则最大值和最小值均在边界取得。若最大值在内部取得，则 $u$ 为常数。
\end{theorem}
\begin{proof}
设 $u(x_0)=M=\max_{\overline\Omega}u$ 且 $x_0\in\Omega$。平均值公式与 $u\le M$ 说明每个充分小球上 $u=M$，否则连续性会使平均严格小于 $M$。因此最大值点集在 $\Omega$ 中既闭又开，连通性使其等于 $\Omega$。若没有内部最大值，则最大值自然在边界；最小值对 $-u$ 应用同一论证。
\end{proof}
对两个调和延拓之差应用定理得
\[
 \|u_1-u_2\|_{C(\overline\Omega)}\le\|g_1-g_2\|_{C(\partial\Omega)}.
\]
这给出唯一性与稳定性，但存在性仍需另行构造。

\subsection{从调和方程推广到泊松方程的比较}
若 $v\in C^2(\Omega)\cap C(\overline\Omega)$ 满足 $\Delta v\ge0$，称为经典\emph{次调和函数（subharmonic function）}，它同样满足最大值在边界的弱结论。证明可用 $v+\varepsilon|x|^2$：其拉普拉斯算子严格为正，不可能在内部取得最大值，最后令 $\varepsilon\downarrow0$。因此，若 $-\Delta u_1\le-\Delta u_2$ 且边界上 $u_1\le u_2$，对 $v=u_1-u_2$ 使用该结论即得内部 $u_1\le u_2$。

\begin{example}[正源项产生非负解]
在球 $B_R(0)\subset\mathbb R^n$ 上，$-\Delta u=1$、$u|_{\partial B_R}=0$ 的解为
\[
 u(x)=\frac{R^2-|x|^2}{2n}.
\]
因为 $\Delta|x|^2=2n$，代入便可核查。它在内部为正并在球心取得最大值，这不违反调和函数的最大值原理，因为此处有源项。一般 $-\Delta u=f\ge0$、零边界的问题，由比较原理也有 $u\ge0$。若 $|f|\le F$，以 $\pm F(R^2-|x|^2)/(2n)$ 比较，还得到 $|u(x)|\le F(R^2-|x|^2)/(2n)$。
\end{example}

\section{泊松积分公式与圆盘边值问题}
先从分离变量看单位圆盘的结构。在极坐标中，角频率 $k\ge1$ 的径向方程是 $r^2R''+rR'-k^2R=0$，解为 $Ar^k+Br^{-k}$；零频率时为 $A+B\log r$。若包含圆心并要求有界，就必须舍去 $r^{-k}$ 和 $\log r$。因此，边界傅里叶展开中的常数项原样保留，第 $k$ 个正弦或余弦项乘以 $r^k$。这解释了高频边界起伏向内部快速减弱的现象，也解释了为什么圆环区域不能直接舍去负幂和对数项。

单位圆盘上连续边值 $g(\theta)$ 的调和延拓为
\begin{equation}
 u(r,\theta)=\frac1{2\pi}\int_0^{2\pi}
 \frac{1-r^2}{1-2r\cos(\theta-\varphi)+r^2}g(\varphi)\,d\varphi,
 \qquad 0\le r<1.
 \label{eq:poisson-disk}
\end{equation}
核的傅里叶展开为 $1+2\sum_{k\ge1}r^k\cos(k(\theta-\varphi))$，由几何级数求和得到。对 $r<1$ 可逐项微分，且每个 $r^k\cos(k\theta)$、$r^k\sin(k\theta)$ 都调和。核非负、积分为 $2\pi$，当 $r\uparrow1$ 时质量集中在 $\varphi=\theta$ 附近；利用 $g$ 的一致连续性便可证明边界一致收敛。

对于 $n\ge2$ 的球 $B_R(0)$，对应\emph{泊松核（Poisson kernel）}为
\[
 P(x,y)=\frac{R^2-|x|^2}{\sigma_{n-1}R|x-y|^n},\qquad |x|<R,\quad |y|=R,
\]
解为 $\int_{\partial B_R}P(x,y)g(y)\,dS_y$。它将在第\ref{chap:green}章中与法向导数表示联系起来。

\section{哈纳克不等式与刘维尔定理}
\begin{proposition}[哈纳克不等式的局部形式]
若 $u\ge0$ 在 $B_{2R}(x_0)$ 调和，则对 $x\in B_R(x_0)$ 有 $C_n^{-1}u(x_0)\le u(x)\le C_nu(x_0)$，其中 $C_n$ 只依赖维数。
\end{proposition}
\begin{proof}
在半径 $3R/2$ 的球面使用泊松表示。对 $|x-x_0|\le R$，分子介于固定正数乘 $R^2$ 之间，$|x-y|$ 介于 $R/2$ 与 $5R/2$ 之间，所以核与球心处的常数核上下可比。乘以非负边值并积分即得。
\end{proof}
沿紧子集中的有限球链迭代可比较正调和函数在不同点的值，这称\emph{哈纳克不等式（Harnack inequality）}。

\begin{theorem}[刘维尔定理]
定义在整个 $\mathbb R^n$ 上且有界的调和函数必为常数。
\end{theorem}
\begin{proof}
在任意球 $B_R(x_0)$ 的泊松公式中对 $x$ 求导并取 $x=x_0$，可得 $|\nabla u(x_0)|\le C_nR^{-1}\|u\|_\infty$。令 $R\to\infty$ 得梯度为零，再用连通性得常数。
\end{proof}

\section{例题与习题}
\begin{example}
求单位圆盘上边值 $g(\theta)=1+2\cos2\theta-\sin3\theta$ 的调和延拓。
\end{example}
\begin{solve}
每个角频率 $k$ 对应原点正则的径向因子 $r^k$，所以
\[
 u(r,\theta)=1+2r^2\cos2\theta-r^3\sin3\theta.
\]
逐项验证调和与边值，再由最大值原理保证唯一性。取球心值得 $u(0)=1$，恰等于边界平均值。
\end{solve}
\begin{enumerate}
 \item 求二维径向调和函数，解释包含原点时对数项为什么必须消失。
 \item 证明 $x_1$ 是全空间调和但非常数的函数，指出它为何不违背刘维尔定理。
 \item 若连续边值满足 $m\le g\le M$，用泊松核证明解也满足相同界。
\end{enumerate}

